\scanpage{86}
% archive foliation: 87

limite\inserted{pour les $G(n)$}, on voit que tout groupe de BT $G$ sur $S$ \deleted{est extension d'un groupe de BT $\unsure{G(}\quad G^o\overset{\mathrm{dfn}}{=}\lim G(n)$} admet une filtration par des sous-groupes de BT $G^o=\lim G(n)^o$ (plusgrand sous-groupe de BT à fibres connexes) et\inserted{(si car. $k(s)=p$) par} $G^{\mathrm{tm}}=\lim G(n)^{\mathrm{tm}}$ (plus grand sous-groupe de BT \deleted{à fibres} toroidal), canonique et \unsure{fonctorielle|foncturielle} en $G$. En particulier, il est extension d'un groupe de BT ind-étale $G/G^o$ par le groupe de BT à fibres connexes $G^o$, qui peut s'interpréter (si $k(s)$ de car. $p$ \deleted{)comme} et l'idéal maximal de $\underline{O}_{S,s}$ nilpotent, p.ex. $S$ noethérien) comme un groupe de Lie formel, et il pourra même s'interpréter comme extension triple d'un\deleted{e} ind-étale par un connexe unipotent par un toroidal.

Sur une base non ponctuelle (p.ex. un trait de car. $p$ )\deleted{0)}, il n'est plus vrai en général qu'un groupe de BT $G$ \deleted{admette} soit extension d'un groupe de BT ind-étale par un groupe de BT à fibres connexes. Il suffit encore de prendre le groupe \deleted{$\unsure{G(\infty)}$} $A(\infty)$, $A$ courbe elliptique modulaire en car. $p$ : l'ennui provient encore des points où $A_s$ est d'invariant de Hasse nul. On a la

\uline{Proposition}.\quad Conditions équivalentes sur le groupe de BT $G$ sur $S$:

(i)\quad $G$ est extension d'un groupe de BT ind-étale par un groupe de BT à fibres connexes.

(ii)\quad Pour tout $n$, $G(n)$ est extension d'un groupe fini étale par un groupe fini radiciel.

(iii)\quad Kif-kif avec $n=1$.

(iv)\quad La fonction $s\mapsto\operatorname{rang\ Sép}G_s=\operatorname{card}G(\bar{s})$ est localement constante.

Résulte facilement du résultat suivant (qui figure dans EGA IV ?):

\uline{Lemme}.

Si $\unsure{X}\longrightarrow S$ est un morphisme fini localement libre, alors $X\longrightarrow S$ se facto-
